% !TEX root = CoeneDylanBP.tex
%%=============================================================================
%% Verfijning via skeletdetectie
%%=============================================================================

\section{Pose Estimation}
\label{sec:pose}

Indien de baseline uit Fase~3 een voldoende proof-of-concept oplevert, wordt het zwaartepunt vervangen door een nauwkeuriger referentiepunt via \textit{Pose Estimation}. Het doel is de systematische fout door bounding box-variatie te elimineren.

Om de positie langs de korte zijde van de trampoline te bepalen, wordt louter in het geïdentificeerde \textit{landingsframe} (bepaald in Fase~\ref{sec:tof}) een Pose Estimation model ingezet. Dit detecteert specifieke lichaamspunten (\textit{keypoints}) van de springer. Om de landingslocatie te specificeren wordt specifiek gezocht naar de enkels. Het berekend middelpunt tussen deze twee keypoints dient als een nauwkeurig referentiepunt voor de landing van de voeten. Dit elimineert de foutgevoeligheid veroorzaakt door vari{\"e}rende armhouding of lichaamsrotatie van de springer die zich in de volledige bounding box voordoen.

Twee kandidaatmodellen worden ge{\"e}valueerd:
\begin{itemize}
    \item \textbf{ViTPose}~\autocite{Xu2022}: een Vision Transformer-gebaseerd model met state-of-the-art prestaties op \textit{keypoint detection benchmarks}.
    \item \textbf{MediaPipe Pose}~\autocite{Lugaresi2019}: een lichter, real-time alternatief van Google dat geschikt is voor deployment op beperkte hardware.
\end{itemize}

Dit berekende middelpunt in pixelco{\"o}rdinaten wordt vervolgens met behulp van de eerder vermelde homografiematrix $H$ getransformeerd naar het co{\"o}rdinaat in het stelsel van de trampoline. Op basis van deze transformatie wordt bepaald op welke HD-zone de landing wordt geprojecteerd binnen het stelsel van de trampoline. Zo wordt de horizontale verplaatsing nauwgezet vastgesteld.

Doordat het referentiepunt gefixeerd is en niet varieert met armbeweging, wordt een lagere systematische fout verwacht ten opzichte van de bounding box-zwaartepuntbenadering. De vergelijking van beide benaderingen vormt een deelresultaat van het onderzoek.

De validatie verloopt op identieke wijze als in Fase~\ref{sec:yolo} (accuracy en gemiddelde afwijking in meters), maar met de Pose Estimation pipeline als detector.
